\section{Temporal Poisson brackets in dynamical cohomology}
\label{sec:temporal-resonance}
This section proves an unrestricted finite-Fourier theorem and a complete
classification of quadratic temporal cancellations. The conclusions concern
classical dynamical cohomology. Neither microscopic BGS universality nor
priority over all external literature is asserted.

\subsection{Definitions and the Fourier criterion}
Let $A\in\mathrm{SL}(2,\mathbb Z)$, $|\operatorname{tr}A|>2$, let
$S=A^{\mathsf T}$ act on $\mathbb T^2$, and write $Cf=f\circ S$.
For a nonconstant real trigonometric polynomial $g$, subtract its mean and put
\[
 V=\{v:\widehat g(v)\ne0\}\subset\mathbb Z^2\setminus\{0\},
 \qquad h_r=\frac{\{g,C^rg\}}{4\pi}\quad(r\ge1).
\]
Here $V=-V$, but Fourier coefficients need not be real. Our Poisson bracket is
$\{f,q\}=\partial_1f\,\partial_2q-\partial_2f\,\partial_1q$ and
$e_v(x)=\exp(2\pi i v\cdot x)$. Consequently
\begin{equation}\label{eq:temporal-coeff}
 \widehat h_r(k)=-\pi\sum_{u+A^rv=k}
       \det(u,A^rv)\widehat g(u)\widehat g(v).
\end{equation}
In particular $\widehat h_r(0)=0$.

For a mean-zero trigonometric polynomial $f$, let
\[
 H_f(\mathcal O)=\sum_{k\in\mathcal O}\widehat f(k),\qquad
 \mathcal V_A(f)=\sum_{\mathcal O}|H_f(\mathcal O)|^2,
\]
where $\mathcal O$ runs over oriented $A$-orbits, not orbits modulo sign.
These are finite sums. Expanding the squared norm and using Fourier
orthogonality gives
\[
 \lim_{n\to\infty}\frac1n\left\|\sum_{j=0}^{n-1}C^jf\right\|_2^2
 =\sum_{q\in\mathbb Z}\langle C^qf,f\rangle
 =\mathcal V_A(f).
\]
Only finitely many correlations are nonzero: a nonzero integer frequency
cannot return to itself under a nonzero power of $A$.
Moreover $\mathcal V_A(f)=0$ if and only if $f=(C-I)F$ for a trigonometric
polynomial $F$, if and only if it is such a coboundary with $F\in L^2$.
Indeed, on each orbit the Fourier equation is a difference equation on
$\mathbb Z$ with finitely supported forcing. Square-summability forces both
constant tails to vanish, hence the forcing has zero sum. Conversely a
zero-sum finite forcing has a finite cumulative-sum solution. This also proves
the criterion without invoking a regularity theorem for transfer functions.

\subsection{A support lemma uniform in the temporal lag}
\begin{lemma}[Isolated highest-lag frequency]\label{lem:temporal-support}
For every finite nonempty $V=-V\subset\mathbb Z^2\setminus\{0\}$ and every
$r\ge1$, there are $u,v\in V$ and $k=u+A^rv$ such that this ordered
decomposition is unique, $D=\det(u,A^rv)\ne0$, and
\[
 k\notin\bigcup_{j=1}^{r-1}(V+A^jV),\qquad
 A^qk\notin\bigcup_{j=1}^{r}(V+A^jV)\quad(q\in\mathbb Z\setminus\{0\}).
\]
\end{lemma}
\begin{proof}
Set $s=\operatorname{sign}(\operatorname{tr}A)$, $B=sA$,
$t=|\operatorname{tr}A|\ge3$, and
\[
 \lambda=\frac{t+\sqrt{t^2-4}}2,\qquad \Lambda=\lambda^r.
\]
We use $\lambda-\lambda^{-1}>2$ and
$\lambda^2-2\lambda-1=(t-2)\lambda-2>0$.
Choose independent real eigen-covectors $\xi,\eta$ of $B$, with factors
$\lambda,\lambda^{-1}$. Neither vanishes on a nonzero integer vector:
$t^2-4$ lies strictly between $(t-1)^2$ and $t^2$, so the eigendirections
are irrational. Set $X=\max_V|\xi|>0$ and $Y=\max_V|\eta|>0$.
Choose $v_X,u_Y\in V$ attaining $\xi(v_X)=X$, $\eta(u_Y)=Y$, and put
$a=\xi(u_Y)/X$, $b=\eta(v_X)/Y$. Choose $\sigma\in\{1,-1\}$ with
$\sigma(a+b)\ge0$ and define
\[
 L(z)=\xi(z)/X+\Lambda\sigma\eta(z)/Y.
\]
The trial points $\sigma u_Y\in V$ and $B^rv_X\in B^rV$ show that the sum
of the two separate maxima of $L$ is at least $2\Lambda+|a+b|$.
Select the lexicographic maximizers of $(L,\xi)$ on $V$ and $B^rV$.
The two forms are independent, so both maximizers are unique. Write them as
$u$ and $B^rw$, put $v=s^rw$, and let $k=u+B^rw=u+A^rv$. Then
\begin{equation}\label{eq:temporal-large-L}
 L(k)\ge2\Lambda.
\end{equation}
An alternative ordered decomposition of $k$ would sum two lexicographically
bounded coordinate pairs to the sum of their maxima. Both pairs must attain
their separate maxima. This proves uniqueness, including ties in $L$.

Every point of $V+A^jV$ satisfies
$|\xi(z)|\le(1+\lambda^j)X$ and
$|\eta(z)|\le(1+\lambda^{-j})Y$. In particular $k$ satisfies these bounds
with $j=r$. If $A^qk\in V+A^jV$ for $q\ge1$, $j\le r$, the first bound
implies $|\xi(k)|\le(1+\Lambda)X/\lambda$. Using $k$'s own second bound,
\[
 L(k)\le\Lambda+\Lambda/\lambda+1+1/\lambda<2\Lambda.
\]
The gap is minimized at $r=1$ and is then $\lambda-\lambda^{-1}-2>0$.
For $q\le-1$ and $j\ge1$ the second bound instead implies
$|\eta(k)|\le(1+\lambda^{-1})Y/\lambda$. Thus
\[
 L(k)\le1+\Lambda+\Lambda/\lambda+\Lambda/\lambda^2<2\Lambda.
\]
Here too the gap is minimized at $r=1$ and equals
$\lambda-\lambda^{-1}-2$. This excludes all nonzero iterates, without a
finite orbit cutoff and for either sign of the trace.
For $1\le j<r$, membership of $k$ itself would give
\[
 L(k)\le1+\Lambda+\lambda^j+\lambda^{r-j}
       \le1+\Lambda+\lambda+\lambda^{r-1}<2\Lambda.
\]
The middle inequality follows from
$(\lambda^{j-1}-1)(\lambda^{r-j-1}-1)\ge0$.
The last gap is at least $\lambda^2-2\lambda-1>0$ since $r\ge2$.

It remains to exclude $D=0$. If $u=cB^rw$, put $R=|c|$,
$a_0=|\xi(w)|/X$, $b_0=|\eta(w)|/Y$.
Membership of $u,w$ in $V$ gives $a_0,b_0\le1$,
$R\Lambda a_0\le1$, $R\Lambda^{-1}b_0\le1$. For $R>0$,
\[
 |L(k)|\le(1+R)\{\min(\Lambda,R^{-1})+\min(1,\Lambda/R)\}
          \le\Lambda+2+\Lambda^{-1}<2\Lambda.
\]
To verify the middle inequality, divide into $R\le\Lambda^{-1}$,
$\Lambda^{-1}\le R\le\Lambda$, and $R\ge\Lambda$.
The expressions are respectively $(1+R)(\Lambda+1)$,
$2+R+R^{-1}$, and $(1+R^{-1})(\Lambda+1)$, maximized at the indicated
boundary points. For $R=0$ the bound $\Lambda+1$ suffices.
All these conclusions contradict \eqref{eq:temporal-large-L}.
\end{proof}

\subsection{Independence and a covariance determinant bound}
\begin{theorem}\label{thm:temporal-independent}
The classes $[h_1],[h_2],\ldots$ are linearly independent modulo $L^2$
$C$-coboundaries. If $a_m\ne0$ and $m$ is the highest lag in a finite real
combination, the witness in Lemma~\ref{lem:temporal-support} gives
\begin{equation}\label{eq:temporal-lower}
 \mathcal V_A\left(\sum_{r=1}^m a_rh_r\right)
 \ge2\pi^2a_m^2D_m^2
       |\widehat g(u_m)\widehat g(v_m)|^2>0.
\end{equation}
\end{theorem}
\begin{proof}
On the orbit of $k_m$ the entire combination has just one coefficient,
$-\pi a_mD_m\widehat g(u_m)\widehat g(v_m)$, by the support lemma and
\eqref{eq:temporal-coeff}. It is nonzero. The orbit of $-k_m$ is distinct,
since $A^qk_m=-k_m$ would require an eigenvalue $-1$ of $A^q$, impossible
for $q\ne0$, while $q=0$ would force $k_m=0$. Reality supplies the conjugate
coefficient on that second orbit, proving the factor two and the theorem.
\end{proof}

Let $G_m$ be the covariance matrix
$(G_m)_{ij}=\sum_{\mathcal O}H_{h_i}(\mathcal O)
\overline{H_{h_j}(\mathcal O)}$. It is real positive definite. More precisely,
\begin{equation}\label{eq:temporal-det}
 \det G_m\ge(2\pi^2)^m
       \prod_{r=1}^m D_r^2|\widehat g(u_r)\widehat g(v_r)|^2.
\end{equation}
Indeed the chosen orbits for distinct lags, including their negatives, are
pairwise distinct. Their orbit-sum rows have zeros in columns $j<r$ and
nonzero diagonal entries $z_r=-\pi D_r\widehat g(u_r)\widehat g(v_r)$.
The Cauchy--Binet formula includes $2^m$ triangular minors obtained by choosing
one member of each conjugate pair; each has squared modulus
$\prod_r|z_r|^2$. All other terms are nonnegative.

\subsection{A complete classification of quadratic temporal cancellations}
\begin{theorem}\label{thm:temporal-cancellation}
Given real $c_{ij}$, $0\le i<j\le m$, define
\[
 H=\frac1{4\pi}\sum_{0\le i<j\le m}c_{ij}\{C^ig,C^jg\},
 \qquad a_r=\sum_{i=0}^{m-r}c_{i,i+r}.
\]
Then $H$ is an $L^2$ coboundary if and only if $a_1=\cdots=a_m=0$.
The coefficient-to-cohomology map has rank $m$ and kernel dimension
$m(m-1)/2$.
\end{theorem}
\begin{proof}
Since $\det S=1$, the Poisson bracket commutes with $C$, and
$\{C^ig,C^{i+r}g\}/(4\pi)=C^ih_r$. Therefore
\[
 H=\sum_{r=1}^m a_rh_r+(C-I)F,\qquad
 F=\sum_{r=1}^m\sum_{i=0}^{m-r}c_{i,i+r}
                      \sum_{\ell=0}^{i-1}C^\ell h_r.
\]
An empty sum is zero. The independence theorem proves necessity; this formula
proves sufficiency with a polynomial transfer. The map $c\mapsto a$ is
surjective onto $\mathbb R^m$, which proves the rank and nullity. A kernel
basis is given by $E_{i,i+r}-E_{0,r}$, $1\le i\le m-r$.
\end{proof}

\begin{corollary}[Finite temporal filters]\label{cor:temporal-filters}
For real polynomials $P,Q$, with $P\ne0$,
\[
 \begin{gathered}
 \{P(C)g,Q(C)g\}/(4\pi)\text{ is a coboundary}\\
 \Longleftrightarrow P(z^{-1})Q(z)=P(z)Q(z^{-1})\\
 \Longleftrightarrow Q/P\in\mathbb R(z+z^{-1}).
 \end{gathered}
\]
\end{corollary}
\begin{proof}
Extend coefficient lists by zeros. The coefficient of $z^r$ in the difference
of the two Laurent polynomials is
$\sum_i(p_iq_{i+r}-q_ip_{i+r})$, exactly the diagonal sum in the preceding
theorem. The constant coefficient is zero and negative coefficients are the
negatives of the positive ones. For the last equivalence put $w=z+z^{-1}$.
The extension $\mathbb R(z)/\mathbb R(w)$ has basis $1,z$, and its involution
sends $z$ to $w-z$. An invariant element $a(w)+zb(w)$ must have $b=0$.
\end{proof}

The criterion permits nonzero brackets with zero variance: for $P=z$ and
$Q=1+z^2$ it gives
\[
 \{Cg,g+C^2g\}/(4\pi)=(C-I)h_1.
\]
Thus nonvanishing of a Poisson bracket must not be confused with a nonzero
cohomology class.

\begin{corollary}[Autocorrelation determines the second-order class]
Write $h(f)=\{f,Cf\}/(4\pi)$. For two real finite filters,
\[
 [h(P(C)g)]=[h(Q(C)g)]
 \quad\Longleftrightarrow\quad
 P(z)P(z^{-1})=Q(z)Q(z^{-1}).
\]
\end{corollary}
\begin{proof}
Let $b_j(P)=\sum_i p_ip_{i+j}$ for $j\ge0$. Applying the cancellation
formula to $P$ and $zP$ gives the coefficient
$b_{r-1}(P)-b_{r+1}(P)$ at $[h_r]$. Equality of classes therefore makes the
difference of the two autocorrelation sequences constant separately on each
parity. Both sequences vanish eventually, so their difference is zero.
The converse follows from the same formula.
\end{proof}
For example $P=1+5z+6z^2$ and $Q=2+7z+3z^2$ are not related by a sign,
shift or reversal, but both have autocorrelations $(62,35,6)$ and give the
same class $56[h_1]+35[h_2]+6[h_3]$. This is equality of a second-order
class, not equality of full quantum spectra or of all perturbative orders.

\subsection{Finite bandwidth and loss of coercivity}
The case $r=1$ proves $\Gamma_A(g):=\mathcal V_A(h_1)>0$ without a
bandwidth or phase restriction. If $M=|V|$ and
$|\widehat g(v)|\ge\alpha\|g\|_2/\sqrt M$ on $V$, then the integer
$D_1\ne0$ in \eqref{eq:temporal-lower} yields
\[
 \Gamma_A(g)\ge2\pi^2\alpha^4M^{-2}\|g\|_2^4.
\]
The amplitude hypothesis cannot be dropped. Let $\mathcal P_R$ be the real
mean-zero Fourier polynomials of max-norm bandwidth at most $R\ge1$ and
$c_R(A)=\inf_{0\ne g\in\mathcal P_R}\Gamma_A(g)/\|g\|_2^4$.
Compactness of the finite-dimensional unit sphere gives $c_R(A)>0$.
Nevertheless, for every $s>1$ there are constants depending on $A,s$ with
\begin{equation}\label{eq:temporal-gevrey}
 0<c_R(A)\le C_s\exp(-c_sR^{1/s}).
\end{equation}
To see this, choose a small box $U$ around a point outside $Sx=\pm x$ so
that $U\cup(-U)$ is disjoint from its inverse image under $S$.
Take a nonzero real even Gevrey-$s$ bump supported there, and subtract its
mean to obtain $b$. Then $\{b,Cb\}=0$, since the derivative supports are
disjoint. Such bumps are obtained from
$\exp(-(1-x^2)^{-1/(s-1)})$ on $(-1,1)$, extended by zero, by products,
translation and rescaling. Their derivatives obey $C L^n(n!)^s$.
For completeness, this derivative estimate follows from Cauchy's formula on
a complex disk of radius proportional to the distance $d$ to the endpoint:
the bound is $n! (c d)^{-n}\exp(-c'd^{-1/(s-1)})$. Maximizing in $d$
gives $C L^n(n!)^s$; all derivatives tend to zero at the endpoint.
Integration by parts and optimization in $n$ imply
$|\widehat b(v)|\le C\exp(-c|v|^{1/s})$.
For $b_R=P_Rb$ and $M_1(f)=\sum_v|v||\widehat f(v)|$, this gives
$M_1(b-b_R)\le C\exp(-cR^{1/s})$ after absorbing polynomial factors.
The determinant bound in the Fourier bracket implies
\[
 \left\|\widehat{\{b_R,Cb_R\}/(4\pi)}\right\|_1
 \le\pi\|A\|_{\mathrm{op}}
       \{M_1(b_R)+M_1(b)\}M_1(b_R-b).
\]
Orbit sums satisfy $\mathcal V_A(f)\le\|\widehat f\|_1^2$.
Also $\|b_R\|_2$ stays bounded below. This proves
\eqref{eq:temporal-gevrey}, using strict positivity for every nonzero $b_R$.
No nondegeneracy theorem for arbitrary infinite analytic Fourier support is
inferred from these finite-band statements.

\subsection{Arbitrarily high-order collective soft modes}
The following result distinguishes finite-dimensional nondegeneracy from
coercivity uniform in the number of temporal lags.
Let $t=|\operatorname{tr}A|$, $\delta=\sqrt{t^2-4}$,
$\lambda=(t+\delta)/2$, and $g_0=a_0e_v+\overline{a_0}e_{-v}$,
where $v\ne0$ is integral and $a_0\ne0$.

\begin{lemma}[Exact single-mode orthogonality]\label{lem:single-lag-orth}
For $D_j=\det(v,A^jv)$, the lag classes generated by $g_0$ have covariance
\[
 \langle[h_j(g_0)],[h_k(g_0)]\rangle
       =4\pi^2|a_0|^4D_j^2\,\mathbf1_{j=k},\qquad
 D_j^2=D_1^2\frac{(\lambda^j-\lambda^{-j})^2}{\delta^2}.
\]
\end{lemma}
\begin{proof}
The four Fourier frequencies of $h_j(g_0)$ are
$\pm(v+A^jv)$ and $\pm(v-A^jv)$; each coefficient has modulus
$\pi|D_j||a_0|^2$. The invariant $q_A(w)=\det(w,Aw)$ satisfies
\[
 q_A(v\pm A^jv)=q_A(v)(2\pm\operatorname{tr}A^j).
\]
The numbers $|\operatorname{tr}A^j|=\lambda^j+\lambda^{-j}$ strictly
increase and exceed two. These invariants distinguish the two types and all
different positive lags, including negative trace. Opposite frequencies
cannot share an oriented orbit. The covariance formula follows.
Cayley--Hamilton gives $D_{j+1}=\operatorname{tr}(A)D_j-D_{j-1}$ with
$D_0=0$, which proves the squared closed form.
\end{proof}

\begin{theorem}[Soft modes of every even differential order]
\label{thm:temporal-soft}
Fix an integer $q\ge1$ and set
\[
 P_q(z)=(z^2-tz+1)^q,\qquad g_q=P_q(C)g_0,\qquad d=2q.
\]
For any fixed $R\ge d$, consider the $n$ cohomology vectors
$\lambda^{-r}[h_r(g_q)]$, $r=R+1,\ldots,R+n$, and their covariance matrix
$\widetilde G_{R,n}$. There are positive constants, independent of $n$, such
that
\[
 c\,n^{-4q}\le\lambda_{\min}(\widetilde G_{R,n})\le C\,n^{-4q}.
\]
Thus the finite matrices are all positive definite but have collective
almost-null directions of arbitrarily high algebraic order.
\end{theorem}
\begin{proof}
Write $P_q(z)=\sum_{i=0}^dp_iz^i$ and
$b_k=\sum_i p_ip_{i+k}$, with $b_{-k}=b_k$ and $b_k=0$ outside $[-d,d]$.
For $r>d$, symplectic transport and reduction modulo coboundaries give the
exact identity
\[
 [h_r(g_q)]=\sum_{k=-d}^d b_k[h_{r+k}(g_0)].
\]
Let $w_k=b_k\lambda^k$ and
$C_0=4\pi^2|a_0|^4D_1^2/\delta^2$. By the preceding lemma, for a real
coefficient vector $x=(x_r)$ the squared norm of its combination is
\[
 x^{\mathsf T}\widetilde G_{R,n}x
 =C_0\sum_j(1-\lambda^{-2j})^2
            \left|\sum_{r=R+1}^{R+n}x_rw_{j-r}\right|^2.
\]
All relevant $j$ satisfy $j\ge R+1-d\ge1$. Hence, in the order of quadratic
forms,
\begin{equation}\label{eq:soft-sandwich}
 C_0(1-\lambda^{-2(R+1-d)})^2 T_n(F_q)
 \ \le\ \widetilde G_{R,n}\ \le\ C_0T_n(F_q),
\end{equation}
where $T_n(F_q)$ is the Toeplitz matrix with symbol
$F_q(\theta)=|\sum_kw_ke^{ik\theta}|^2$ and Haar normalization
$d\theta/(2\pi)$.
The autocorrelation identity gives
\[
 \sum_kw_ke^{ik\theta}=P_q(z)P_q(z^{-1}),\qquad z=\lambda e^{i\theta}.
\]
Since $(z^2-tz+1)(z^{-2}-tz^{-1}+1)=(z+z^{-1}-t)^2$, putting
$u(\theta)=2-2\cos\theta$ yields the exact factorization
\begin{equation}\label{eq:soft-symbol}
 \boxed{F_q(\theta)=u(\theta)^{2q}(\delta^2+u(\theta))^{2q}.}
\end{equation}
Indeed $|z+z^{-1}-t|^2=t^2(\cos\theta-1)^2+\delta^2\sin^2\theta
=u(\delta^2+u)$.

It remains to bound the smallest eigenvalue. Since $0\le u\le4$, the factor
$(\delta^2+u)^{2q}$ is bounded above and below by positive constants.
For $p(\theta)=\sum_{j=1}^nx_je^{ij\theta}$ with $\|x\|_2=1$,
Jensen's inequality and the Dirichlet tridiagonal Laplacian give
\[
 \int u^{2q}|p|^2\frac{d\theta}{2\pi}
 \ge\left(\int u|p|^2\frac{d\theta}{2\pi}\right)^{2q}
 \ge\left(4\sin^2\frac{\pi}{2(n+1)}\right)^{2q}.
\]
For an upper bound choose $x_j=\phi(j/(n+1))$, where
$\phi(x)=x^{2q}(1-x)^{2q}$ on $[0,1]$ and is zero elsewhere.
This extension has bounded weak derivative of order $2q$.
Its finite differences of order $2q$ are bounded by $C_q(n+1)^{-2q}$,
including across the endpoints. Parseval identifies the last integral,
before normalization, with the sum of squares of these differences.
That sum is at most $C_q n^{1-4q}$, while $\sum_j|x_j|^2\ge c_qn$
for all sufficiently large $n$. This proves the matching $O(n^{-4q})$
Rayleigh bound; finitely many smaller $n$ are absorbed in the constants.
Combining with \eqref{eq:soft-sandwich} proves the theorem for the actual
finite-lag covariance matrices, not only for a limiting Toeplitz model.
\end{proof}

The normalization $\lambda^{-r}$ removes the natural exponential growth of
individual brackets. The theorem does not assert an exact null vector at any
finite $n$, a random-matrix law, or a uniform-in-time quantum correspondence.
It identifies a quantified obstruction to replacing finite-lag positivity by
a lag-uniform lower bound.

\subsection{A finite quantum consequence and its time boundary}
Use the physical Weyl lift
\[
 (Q_N(a,b)u)(x)=e^{2\pi ia(x-b/2)/N}u(x-b\bmod N),\qquad
 Q_N(v)Q_N(w)=e^{\pi i\det(v,w)/N}Q_N(v+w).
\]
Write $\operatorname{Op}_N(f)=\sum_v\widehat f(v)Q_N(v)$ and, for real
$a_r$ with highest nonzero index $m$, form the Hermitian operator
\[
 T_{n,N}=\sum_{\ell=0}^{n-1}\sum_{r=1}^m a_r\frac{iN}{2}
 [\operatorname{Op}_N(C^\ell g),\operatorname{Op}_N(C^{\ell+r}g)].
\]
These are defined as transported commutators. Identifying them with time
evolution by a quantum cat map additionally requires exact Egorov in the
chosen boundary sector; it is not assumed for every $A$ in a fixed sector.

\begin{proposition}\label{prop:temporal-quantum}
Let $R=\max_{v\in V}\|v\|_\infty$, $a=\|A\|_\infty$. If
\[
                 2R a^{n-1}(1+a^m)<N,
\]
then the witness of Lemma~\ref{lem:temporal-support} at lag $m$ gives
\[
 \frac{\operatorname{Tr}(T_{n,N}^2)}{nN}
 \ge2a_m^2N^2\sin^2(\pi D_m/N)
                   |\widehat g(u_m)\widehat g(v_m)|^2.
\]
It is strictly positive when $N$ does not divide $D_m$.
\end{proposition}
\begin{proof}
All occurring frequencies lie in the open max-norm box of radius $N/2$,
so distinct integer frequencies remain distinct modulo $N$ and their Weyl
operators are orthogonal for the normalized Hilbert--Schmidt inner product.
The support lemma isolates the $2n$ different frequencies
$\pm A^\ell k_m$, $0\le\ell<n$, in the entire transported sum.
Their coefficient moduli are
$|a_m|N|\sin(\pi D_m/N)|\,|\widehat g(u_m)\widehat g(v_m)|$,
since determinants are preserved by $A$. Summing their squared moduli proves
the bound. Hermiticity identifies the squared Hilbert--Schmidt norm with the
trace in the statement.
\end{proof}
For fixed $A,R,m$ the sufficient no-alias condition allows only
$n=O(\log N)$, not the Heisenberg scale $n\asymp N$. This is an observable
variance bound, not a joint-trace estimate for eigenphase correlations.
The distinction is essential: at $N=3$, $A=\left(\begin{smallmatrix}4&1\\3&1\end{smallmatrix}\right)$
and $g=e_{(1,0)}+e_{(-1,0)}$, the classical determinant is $3$, but
$Q_3(A(1,0))=Q_3(1,0)$ and the first quantum commutator vanishes.
The no-alias condition fails in this negative control.
